
\documentclass{article}
\usepackage{colortbl}
\usepackage{makecell}
\usepackage{multirow}
\usepackage{supertabular}

\begin{document}

\newcounter{utterance}

\centering \large Interaction Transcript for game `ifeval', experiment `ifeval', episode 62 with reflexion{-}Qwen3.5{-}9B.
\vspace{24pt}

{ \footnotesize  \setcounter{utterance}{1}
\setlength{\tabcolsep}{0pt}
\begin{supertabular}{c@{$\;$}|p{.15\linewidth}@{}p{.15\linewidth}p{.15\linewidth}p{.15\linewidth}p{.15\linewidth}p{.15\linewidth}}
    \# & \multicolumn{2}{c}{Player} && \multicolumn{2}{c}{Game Master} \\
    \hline

    \theutterance \stepcounter{utterance}  
    & & & \multicolumn{4}{p{0.6\linewidth}}{
        \cellcolor[rgb]{0.9,0.9,0.9}{
            \makecell[{{p{\linewidth}}}]{
                \texttt{\tiny{[P1$\langle$GM]}}
                \texttt{Which of the following is not a fish: salmon or avocado? Answer this easy question first, then expand into an interesting article about salmon and avocado. Use some words in all caps in your response, but those all{-}caps words should appear at most 10 times.} \\
            }
        }
    }
     \\ \\

    \theutterance \stepcounter{utterance}  
    & \multicolumn{4}{p{0.6\linewidth}}{
        \cellcolor[rgb]{0.561,0.737,0.561}{
            \makecell[{{p{\linewidth}}}]{
                \texttt{\tiny{[P1$\rangle$GM]}}
                \texttt{The answer is simple: **AVOCADO** is not a fish.} \\
\\ 
\texttt{While salmon swims in the ocean and is a delicious source of protein, the avocado is a creamy fruit native to Mexico. It is often called a "superfood" for its healthy fats, but it definitely doesn't have gills or fins.} \\
\\ 
\texttt{***} \\
\\ 
\texttt{\#\#\# The Dynamic Duo: Salmon and Avocado} \\
\\ 
\texttt{Despite their biological differences, salmon and avocado have become an iconic pairing in modern cuisine. This combination works because they are both powerhouses of nutrition, offering a perfect balance of macronutrients that fuel the body.} \\
\\ 
\texttt{**The Omega{-}3 Powerhouse**} \\
\texttt{Salmon is renowned for its high content of **OMEGA{-}3** fatty acids. These essential fats are crucial for heart health, brain function, and reducing inflammation. When you eat wild{-}caught salmon, you are also getting astaxanthin, the pigment that gives the fish its vibrant pink color and provides powerful antioxidant benefits.} \\
\\ 
\texttt{**The Healthy Fat Carrier**} \\
\texttt{Enter the avocado. While it is a fruit, it is packed with monounsaturated fats, specifically oleic acid, which is the same healthy fat found in olive oil. The beauty of pairing avocado with salmon lies in absorption. The fats in the avocado help your body absorb the fat{-}soluble vitamins (like D and E) found in the fish, making the meal even more nutritious.} \\
\\ 
\texttt{**A Culinary Symphony**} \\
\texttt{From the classic salmon sashimi topped with a slice of ripe avocado to grilled salmon fillets drizzled with a creamy avocado salsa, the possibilities are endless. The texture contrast is delightful: the firm, flaky flesh of the fish against the buttery, smooth creaminess of the fruit.} \\
\\ 
\texttt{Whether you are looking to boost your energy levels, support your heart, or simply enjoy a meal that tastes incredible, combining these two ingredients is a WIN{-}WIN for both your palate and your health.} \\
            }
        }
    }
    & & \\ \\

    \theutterance \stepcounter{utterance}  
    & & & \multicolumn{2}{p{0.3\linewidth}}{
        \cellcolor[rgb]{0.9,0.9,0.9}{
            \makecell[{{p{\linewidth}}}]{
                \texttt{\tiny{[GM$|$GM]}}
                {'change_case:capital_word_frequency': {'capital_relation': 'at least', 'capital_frequency': 1}}
            }
        }
    }
    & & \\ \\

    \theutterance \stepcounter{utterance}  
    & & & \multicolumn{2}{p{0.3\linewidth}}{
        \cellcolor[rgb]{0.9,0.9,0.9}{
            \makecell[{{p{\linewidth}}}]{
                \texttt{\tiny{[GM$|$GM]}}
                \texttt{game\_result = WIN} \\
            }
        }
    }
    & & \\ \\

\end{supertabular}
}

\end{document}
